\chapter{Calibration and Sensor Models}

In Chapter~34 (Section~\ref{ch34:three-questions}) we met the three core sensors of motion sensing ---
the accelerometer, the gyroscope, and the magnetometer --- and saw, in principle,
what each one measures. The qualifier ``in principle'' matters. A
datasheet promises that an accelerometer reports specific force in
$\mathrm{m/s^2}$ and a magnetometer reports the field in microtesla. Wire one up,
hold it dead still on your desk, and the numbers it gives you will be
\emph{wrong}. The errors are usually small --- a few percent here, a
fraction of a degree there, a slow creep over a minute --- but they corrupt
everything downstream. Feed those readings into an orientation
filter (next chapter) and the errors do not average out; they accumulate, and
your virtual wrist drifts away from your real one.

The fix is \emph{calibration}: building a small mathematical model of how the
sensor lies, measuring the parameters of that model, and then undoing the lie in
software before any other processing touches the data. This chapter is about
that model and that measurement. The chapter is both practical and geometric: the central
method, the ellipsoid fit, treats the sensor errors as squashing and shifting a
sphere, and calibration undoes both.

We will build the standard sensor error model one term at a time, each with a
picture of what it does to your data. We will see why an ideal accelerometer or
magnetometer, swung through all orientations, should trace a perfect sphere ---
and why a real one traces a tilted, off-center ellipsoid instead. We will fit
that ellipsoid and read the calibration off of it. We will estimate gyroscope
bias and watch it drift. We will meet the Allan variance, a plot that shows how a
noisy sensor's errors behave over different time scales. And
throughout, we will be clear about what calibration can and cannot rescue:
no amount of clever fitting recovers individual finger joint angles from wrist
bulk motion.

\section{Why a raw reading is not the truth}
\label{ch35:why-raw-lies}

Start with the cleanest possible mental model of a sensor. You have some
physical quantity --- call it the \emph{true} value $\vec{a}$, a 3-vector for an
accelerometer or magnetometer, a scalar per axis for a rate gyro. The sensor
performs a measurement and hands you back a number $\vec{y}$. In a perfect world
$\vec{y} = \vec{a}$ and we would never have written this chapter. In the real
world there is a \emph{transfer relationship} between them, and
calibration means discovering that relationship and inverting it.

\begin{intuitionbox}
Think of the sensor as a slightly drunk translator. You hand it the truth in one
language (real acceleration in $\mathrm{m/s^2}$) and it hands you back a
translation (counts, or volts, or scaled floats). A \emph{good} translator
gets it right. A drunk one consistently adds a fixed greeting you never asked for
(bias), stretches every sentence by ten percent (scale factor), occasionally
swaps two words (axis misalignment), gets sloppier as the room warms up
(temperature drift), and slurs randomly from word to word (noise). Calibration
is learning the translator's specific quirks well enough to mentally subtract
them --- so that even a drunk translator becomes reliable, as long as it is
\emph{consistently} drunk.
\end{intuitionbox}

The key word there is \emph{consistently}. Calibration only helps with the
parts of the error that are \emph{systematic} --- repeatable, the same today as
yesterday, a deterministic function of the true value. The part that is genuinely
random --- the slurring --- cannot be calibrated away. It can only be
\emph{characterized} (so you know how much to trust a reading) and \emph{filtered}
(averaged down, as in Chapter~20). So our sensor model splits cleanly in two:
\[
  \underbrace{\text{systematic error}}_{\text{calibrate it out}}
  \;+\;
  \underbrace{\text{random error}}_{\text{characterize, then filter}}.
\]
The rest of Section~\ref{ch35:error-model} builds the systematic part term by
term; Sections~\ref{ch35:allan} and onward handle the random part.

\section{The standard sensor error model}
\label{ch35:error-model}

Here is the model we are heading toward. For a 3-axis sensor, the measured
vector $\vec{y}$ relates to the true vector $\vec{a}$ by
\begin{equation}
  \vec{y} \;=\; \mat{S}\,\mat{M}\,\mat{R}\,\vec{a} \;+\; \vec{b} \;+\; \vec{n},
  \label{ch35:full-model}
\end{equation}
where $\vec{b}$ is a constant \emph{bias} (offset), $\mat{S}$ is a diagonal
\emph{scale-factor} matrix, $\mat{M}$ captures \emph{misalignment} and
cross-axis coupling, $\mat{R}$ is the rotation from the sensor's nominal axes to
its true sensing axes, and $\vec{n}$ is random \emph{noise}. The three matrices
in front are almost always lumped together into one general $3\times3$ matrix,
\begin{equation}
  \vec{y} \;=\; \mat{A}\,\vec{a} \;+\; \vec{b} \;+\; \vec{n},
  \qquad \mat{A} = \mat{S}\,\mat{M}\,\mat{R},
  \label{ch35:lumped-model}
\end{equation}
and calibration becomes: \emph{find $\mat{A}$ and $\vec{b}$, then recover}
$\hat{\vec{a}} = \mat{A}^{-1}(\vec{y}-\vec{b})$. The rest of the
chapter illustrates each term, gives methods for measuring the terms, and notes
their limits. We take the terms one at a time.

\subsection{Bias: a finger on the scale}
\label{ch35:bias}

The simplest error is a constant offset. The sensor adds (or subtracts) a fixed
amount regardless of the true value:
\begin{equation}
  y = a + b.
  \label{ch35:bias-scalar}
\end{equation}
A gyroscope sitting perfectly still should read zero rotation rate; if it reads
$0.6^\circ/\mathrm{s}$, that $0.6$ is bias. An accelerometer lying flat should
report exactly $1\,g$ on its vertical axis and $0$ on the others; a few hundred
$\mathrm{mg}$ on a horizontal axis is bias. Geometrically, bias is the simplest
error --- it just \emph{shifts} every reading by the same vector $\vec{b}$
--- but for integrating sensors it is the most dangerous, because a constant
offset integrates into a ramp. We return to that in
Section~\ref{ch35:gyro-bias}.

\begin{figure}[H]
\centering
\begin{tikzpicture}
\begin{axis}[
title style={font=\small},
    width=11.5cm, height=4.6cm,
    xlabel={true value $a$}, ylabel={measured $y$},
    xmin=-1.1, xmax=1.1, ymin=-1.3, ymax=1.7,
    axis lines=left, samples=2, domain=-1:1,
    legend pos=north west, legend cell align=left,
    title={Three systematic errors as distortions of the input--output line},
    every axis plot/.append style={line width=1.0pt},
]
\addplot[gray, dashed] {x}; \addlegendentry{ideal $y=a$}
\addplot[sigblue] {x + 0.35}; \addlegendentry{bias: $y=a+0.35$}
\addplot[specred] {1.4*x}; \addlegendentry{scale: $y=1.4\,a$}
\addplot[quotientorange] {1.4*x + 0.35}; \addlegendentry{both}
\end{axis}
\end{tikzpicture}
\caption{For a single axis, the ideal sensor is the line $y=a$ (gray dashed).
Bias slides the line vertically; scale-factor error tilts it about the origin;
real sensors do both. Calibration is the act of recovering the slope and
intercept so the line can be straightened back to $y=a$.}
\label{ch35:fig-bias-scale}
\end{figure}

\subsection{Scale factor: the wrong gain}
\label{ch35:scale}

The next error multiplies rather than adds. The sensor's gain is off, so a true
$1\,g$ reads as $1.03\,g$ or $0.97\,g$:
\begin{equation}
  y = s\,a, \qquad s \approx 1.
  \label{ch35:scale-scalar}
\end{equation}
Each axis can have its own scale factor, which is why $\mat{S}$ in
\eqref{ch35:full-model} is diagonal, $\mat{S}=\diag(s_x,s_y,s_z)$. Scale error
\emph{tilts} the input--output line of Figure~\ref{ch35:fig-bias-scale} instead
of sliding it. Combined with bias you get the affine relationship $y = s\,a + b$
--- a line with the wrong slope \emph{and} the wrong intercept --- and recovering
that one line per axis is exactly the textbook ``six-position'' accelerometer
calibration we will do by hand in Example~\ref{ch35:ex-sixpos}.

\begin{trapbox}
Bias and scale are easy to confuse when you only look at one orientation. If a
flat accelerometer reads $1.05\,g$, is that $0.05\,g$ of bias or a $5\%$ scale
error? You \emph{cannot tell from one measurement} --- both predict the same
reading at $+1\,g$. The two only separate when you look at \emph{different} input
levels: flip the sensor over so the true input is $-1\,g$. Bias makes both
readings shift the same way ($+1g\!\to\!1.05$, $-1g\!\to\!-0.95$); pure scale
makes them shift oppositely ($+1g\!\to\!1.05$, $-1g\!\to\!-1.05$). This is why
any trustworthy calibration excites the sensor at \emph{multiple} known inputs.
\end{trapbox}

\subsection{Misalignment and cross-axis coupling}
\label{ch35:misalign}

A 3-axis sensor is supposed to have three
\emph{orthogonal} sensing axes lined up perfectly with the chip's package and,
through it, with your board. Two things go wrong. First, the physical sensing
axes are not perfectly orthogonal to each other --- the MEMS structure was etched
with finite precision. Second, the whole triad is rotated slightly relative to
where you \emph{think} it points, because the die was bonded into the package at a
tiny angle and you soldered the package onto the board at another tiny angle.

Both effects mean that motion along the true $x$ axis leaks a little signal into
the reported $y$ and $z$ channels. That leakage is exactly an off-diagonal entry
in a matrix:
\begin{equation}
  \begin{bmatrix} y_x \\ y_y \\ y_z \end{bmatrix}
  =
  \begin{bmatrix}
    1 & m_{xy} & m_{xz} \\
    m_{yx} & 1 & m_{yz} \\
    m_{zx} & m_{zy} & 1
  \end{bmatrix}
  \begin{bmatrix} a_x \\ a_y \\ a_z \end{bmatrix}.
  \label{ch35:misalign-matrix}
\end{equation}
The diagonal is the wanted response; the small off-diagonal terms $m_{ij}$ (a few
percent or a few tenths of a degree, in radians) are the cross-talk. A pure
rotation $\mat{R}$ is one special, structured case of this (orthogonal, $\det=1$);
genuine non-orthogonality of the axes is another. In practice we do not bother to
separate $\mat{S}$, $\mat{M}$, and $\mat{R}$ --- they multiply into the single
matrix $\mat{A}$ of \eqref{ch35:lumped-model}, and the ellipsoid fit recovers
$\mat{A}$ as a whole. The separation only matters when you care \emph{why} the
matrix is what it is.

\begin{connectionbox}
Equation~\eqref{ch35:misalign-matrix} is a concrete 3D linear map (Chapter~42) from the
true physical vector to the reported sensor vector: its diagonal entries set per-axis gain,
its off-diagonals capture cross-axis coupling, and $\mat{A}^{-1}$ is the change of basis
that restores orthonormal coordinates. When we fit and diagonalize the calibration
ellipsoid below, we are computing the eigenvectors and eigenvalues of
$\mat{A}^{\!\top}\mat{A}$ --- the principal axes and stretch factors of a singular value
decomposition measured directly on hardware.
\end{connectionbox}

\subsection{Temperature drift}
\label{ch35:temp}

In practice, $\mat{A}$ and $\vec{b}$ are not constants. They
drift with temperature. MEMS structures are tiny silicon springs, and a spring's
stiffness changes as it warms; the electronics' reference voltages drift too. So
bias and scale are really functions of temperature $T$:
\begin{equation}
  \vec{b}(T) \approx \vec{b}_0 + \vec{k}_b\,(T - T_0), \qquad
  \mat{A}(T) \approx \mat{A}_0 + \mat{K}_A\,(T - T_0),
  \label{ch35:temp-model}
\end{equation}
to first order around a reference temperature $T_0$. Gyro bias is the worst
case: a cheap MEMS gyro can shift its zero-rate output by several
$^\circ/\mathrm{s}$ across the temperature swing from a cold room to a
self-heated, hand-warmed device. A wrist tracker turned on in a cold room and
then warmed by body heat will see its gyro bias \emph{walk} for the first minute
or two of wear --- a real, repeatable effect you can watch on a plot.

\begin{figure}[H]
\centering
\begin{tikzpicture}
\begin{axis}[
title style={font=\small},
    width=11.5cm, height=4.6cm,
    xlabel={time since power-on (s)}, ylabel={gyro zero-rate ($^\circ$/s)},
    xmin=0, xmax=180, ymin=-0.1, ymax=0.95,
    axis lines=left, samples=120, domain=0:180,
    title={Warm-up bias drift: a still gyro's zero slowly settles},
    every axis plot/.append style={line width=1.0pt},
]
\addplot[sigblue] {0.8*(1 - exp(-x/45)) + 0.03*sin(deg(x/6))};
\addplot[gray, dashed] {0.8};
\node[anchor=west, gray] at (axis cs:95,0.86) {\small settled bias};
\end{axis}
\end{tikzpicture}
\caption{A stationary MEMS gyroscope after power-on. The zero-rate output is not
flat --- it ramps as the chip self-heats and the silicon's properties shift, then
settles. Estimating bias in the first few seconds and trusting it for the next
hour is a classic mistake; bias must be re-estimated, ideally continuously.}
\label{ch35:fig-warmup}
\end{figure}

\begin{trapbox}
A common field failure in cheap IMU projects is calibrating once,
cold, on the bench, and baking those numbers into firmware forever. The device
then warms up on the wrist, the bias walks off your stored value, and the
orientation drifts. Either calibrate at the operating temperature, store a
temperature curve as in \eqref{ch35:temp-model}, or --- best for gyros ---
re-estimate bias on the fly whenever the device is detected to be still (see
Section~\ref{ch35:gyro-bias}).
\end{trapbox}

\subsection{Random walk and the noise terms}
\label{ch35:random-walk}

Finally the genuinely random part, $\vec{n}$. It comes in two flavors that
behave very differently under integration, and the Allan variance is the
standard tool for telling them apart.

The first is plain \emph{white noise}: zero-mean, uncorrelated sample-to-sample
jitter on top of every reading. For a gyro this is called \emph{angle random
walk} (ARW), because when you integrate white rate noise to get angle, the angle
performs a random walk --- it wanders away from the truth with a standard
deviation that grows like $\sqrt{t}$ (Chapter~20's averaging law in reverse:
summing $N$ independent steps gives spread $\propto\sqrt{N}$).

In the second, the \emph{bias itself} is not a fixed number but slowly
random-walks over time. This is \emph{bias instability} or \emph{rate random
walk}. It is what makes ``calibrate the bias once'' fail even at constant
temperature --- you nail the bias at $t=0$, and ten minutes later it has
drifted to a different value.

\begin{intuitionbox}
White noise and bias drift pull in opposite directions when you average. Average
\emph{more} samples and white noise shrinks (good) --- but average over a
\emph{longer time} and bias drift grows (bad), because the bias has had longer to
wander. So there is a sweet-spot averaging time: short enough that the bias has
not moved much, long enough that the white noise has been beaten down. The Allan
variance is precisely the plot that reveals this sweet spot, and we draw it in
Section~\ref{ch35:allan}.
\end{intuitionbox}

\section{The big idea: a still sensor should trace a sphere}
\label{ch35:sphere-idea}

The central geometric idea of the chapter applies to any sensor that measures a
vector of \emph{constant magnitude} regardless of orientation. Two of our three
sensors qualify:

\begin{itemize}
\item A \textbf{stationary accelerometer} feels only gravity. Gravity has a fixed
magnitude ($\approx 9.81\,\mathrm{m/s^2}$, or $1\,g$) and points down. As you
rotate the device through every orientation while holding it still, the gravity
vector \emph{in sensor coordinates} sweeps over all directions but always has the
same length. The tip of that vector traces a \textbf{sphere of radius $1\,g$}.
\item A \textbf{magnetometer} sitting in one location feels the local Earth field
(plus any fixed nearby iron). That field has a fixed magnitude and direction in
the room. Rotate the device through every orientation and, again, the field
vector in sensor coordinates traces a \textbf{sphere}, this time of radius equal
to the local field strength.
\end{itemize}

That is the reference truth: \emph{a perfect sensor, swung through all
orientations, returns points that lie exactly on a sphere centered at the
origin}. We do not need to know which orientation produced which point. We do not
need a turntable or a reference instrument. We only need the geometric fact that
the magnitude is constant. This is what makes accelerometer and magnetometer
calibration self-contained: \emph{the calibration target is free.}

\begin{whybox}
Why does constant magnitude buy us a full calibration? Because the error model
\eqref{ch35:lumped-model} maps that sphere to a predictable shape. Apply a matrix
$\mat{A}$ to every point of a sphere and you get an \emph{ellipsoid} (the matrix
stretches different directions by different amounts and rotates the whole thing).
Then add the bias $\vec{b}$ and the ellipsoid \emph{shifts off the origin}. So a
miscalibrated sensor's data does not lie on a centered sphere --- it lies on a
\textbf{tilted, stretched, off-center ellipsoid}. The calibration problem
therefore reduces to: \emph{find the ellipsoid that best fits my data points, then find
the transform that turns that ellipsoid back into the unit sphere.} That transform
\emph{is} $\mat{A}^{-1}$ and $\vec{b}$.
\end{whybox}

\begin{figure}[H]
\centering
\begin{tikzpicture}[font=\small, scale=0.83, every node/.style={transform shape}, >=Latex]
  % Panel A: Raw Tilted, Off-Center Ellipsoid
  \begin{scope}[shift={(0,0)}]
    \node[font=\footnotesize\bfseries, anchor=south] at (2.2, 3.1) {Raw Data: Shifted \& Tilted Ellipsoid};
    \draw[->, gray!60] (-0.3,0) -- (4.4,0) node[right, font=\scriptsize] {$y_x$};
    \draw[->, gray!60] (0,-1.8) -- (0,2.3) node[left, font=\scriptsize] {$y_y$};
    \coordinate (B) at (1.6, 0.6);
    \fill[black] (B) circle (1.5pt) node[below right, font=\footnotesize] {Bias $\vec{b}$};
    \begin{scope}[shift={(B)}, rotate=28]
      \draw[specred, very thick, dashed] (0,0) ellipse (2.0 and 1.0);
      \foreach \ang in {15, 45, 80, 120, 160, 205, 245, 290, 335} {
        \fill[specred] (\ang:2.0 and 1.0) circle (1.6pt);
      }
      \draw[->, thick, specred!80!black] (0,0) -- (2.0, 0) node[midway, above, font=\tiny] {$\lambda_1$};
      \draw[->, thick, specred!80!black] (0,0) -- (0, 1.0) node[midway, left, font=\tiny] {$\lambda_2$};
    \end{scope}
    \node[font=\scriptsize, align=center, text=specred] at (2.2, -2.1) 
      {$\vec{y} = \mat{A}\vec{a}_{\text{true}} + \vec{b}$\\Offset by $\vec{b}$, distorted by $\mat{A}$};
  \end{scope}

  % Step 1 Arrow: Center
  \draw[->, thick, gray!70] (4.6, 0.6) -- node[above, font=\footnotesize] {Center} 
                            node[below, font=\scriptsize] {$\vec{y}' = \vec{y}-\vec{b}$} (6.2, 0.6);

  % Panel B: Centered Ellipsoid
  \begin{scope}[shift={(7.4,0)}]
    \node[font=\footnotesize\bfseries, anchor=south] at (1.0, 3.1) {Centered at Origin};
    \coordinate (O) at (1.0, 0.6);
    \draw[->, gray!60] (O) ++(-2.0,0) -- ++(3.8,0) node[right, font=\scriptsize] {$y'_x$};
    \draw[->, gray!60] (O) ++(0,-1.8) -- ++(0,3.6) node[left, font=\scriptsize] {$y'_y$};
    \fill[black] (O) circle (1.5pt);
    \begin{scope}[shift={(O)}, rotate=28]
      \draw[quotientorange, very thick, dashed] (0,0) ellipse (2.0 and 1.0);
      \foreach \ang in {15, 45, 80, 120, 160, 205, 245, 290, 335} {
        \fill[quotientorange] (\ang:2.0 and 1.0) circle (1.6pt);
      }
    \end{scope}
    \node[font=\scriptsize, align=center, text=quotientorange!80!black] at (1.0, -2.1) 
      {Bias removed;\\scale \& cross-talk remain};
  \end{scope}

  % Step 2 Arrow: Un-deform
  \draw[->, thick, gray!70] (10.6, 0.6) -- node[above, font=\footnotesize] {Un-deform} 
                            node[below, font=\scriptsize] {$\vec{a}_{\text{cal}} = \mat{W}\vec{y}'$} (12.2, 0.6);

  % Panel C: Restored Unit Sphere
  \begin{scope}[shift={(13.4,0)}]
    \node[font=\footnotesize\bfseries, anchor=south] at (1.0, 3.1) {Calibrated Unit Sphere};
    \coordinate (O2) at (1.0, 0.6);
    \draw[->, gray!60] (O2) ++(-1.8,0) -- ++(3.6,0) node[right, font=\scriptsize] {$a_x$};
    \draw[->, gray!60] (O2) ++(0,-1.8) -- ++(0,3.6) node[left, font=\scriptsize] {$a_y$};
    \fill[black] (O2) circle (1.5pt);
    \draw[targetgreen, very thick] (O2) circle (1.5);
    \draw[targetgreen!40, dashed] (O2) ellipse (1.5 and 0.4);
    \foreach \ang in {15, 45, 80, 120, 160, 205, 245, 290, 335} {
      \fill[targetgreen] (O2) ++(\ang:1.5) circle (1.6pt);
    }
    \draw[->, thick, targetgreen!80!black] (O2) -- ++(1.5, 0) node[midway, below, font=\tiny] {$1\,g$};
    \node[font=\scriptsize, align=center, text=targetgreen!60!black] at (1.0, -2.1) 
      {$\|\vec{a}_{\text{cal}}\| = 1.00\,g$\\(Unit Sphere $S^2$)};
  \end{scope}
\end{tikzpicture}
\caption{The geometry of sensor calibration: transforming a distorted ellipsoid into a unit sphere. Left: when an uncalibrated accelerometer or magnetometer is tumbled through all orientations, raw measurements trace a tilted, off-center ellipsoid due to bias $\vec{b}$ and sensitivity/misalignment matrix $\mat{A}$. Middle: least-squares ellipsoid fitting extracts the center $\vec{b}$, translating the ellipsoid to the origin. Right: applying the whitening transformation $\mat{W} = \mat{A}^{-1}$ rotates and scales the principal axes, restoring all measurements onto the perfect unit sphere $\|\vec{a}_{\text{cal}}\|=1\,g$.}
\label{ch35:fig-ellipsoid}
\end{figure}

\section{The ellipsoid fit}
\label{ch35:ellipsoid-fit}

Let us make the picture into arithmetic. Collect a cloud of measured points
$\vec{y}_k$ by slowly tumbling the (stationary, for the accelerometer)
sensor through as many orientations as you can --- the more uniformly you cover
the sphere, the better. We assert that, after we remove bias and undo $\mat{A}$,
each point should have the known constant magnitude $g$ (or field strength $F$).
That is the constraint
\begin{equation}
  \norm{\mat{A}^{-1}(\vec{y}_k - \vec{b})}^2 = g^2 \quad\text{for all } k.
  \label{ch35:constraint}
\end{equation}
A general ellipsoid is the set of points satisfying a quadratic equation,
\begin{equation}
  (\vec{y}-\vec{b})^{\!\top}\,\mat{Q}\,(\vec{y}-\vec{b}) = 1,
  \qquad \mat{Q} = \frac{1}{g^2}\,\mat{A}^{-\top}\mat{A}^{-1},
  \label{ch35:ellipsoid-eq}
\end{equation}
where $\mat{Q}$ is a symmetric positive-definite matrix encoding the ellipsoid's
shape and orientation, and $\vec{b}$ its center. As written, $\vec{b}$ appears
inside a product with $\mat{Q}$, so the equation is not linear in the unknowns. Expand it
and it becomes linear after a change of variables:
\[
  \vec{y}^{\top}\mat{Q}\vec{y} \;-\; 2\,\vec{b}^{\top}\mat{Q}\vec{y}
  \;+\; \vec{b}^{\top}\mat{Q}\vec{b} \;=\; 1
  \quad\Longrightarrow\quad
  \vec{y}^{\top}\tilde{\mat{Q}}\vec{y} \;+\; 2\,\vec{g}^{\top}\vec{y} \;=\; 1,
\]
where we divided through by $1-\vec{b}^{\top}\mat{Q}\vec{b}$ to get
$\tilde{\mat{Q}}$ (a rescaled $\mat{Q}$) and named $\vec{g} = -\tilde{\mat{Q}}\vec{b}$.
Each point $\vec{y}_k$ now gives one equation that is linear in the nine numbers
(six in symmetric $\tilde{\mat{Q}}$, three in $\vec{g}$). With many points we get an
overdetermined linear system and solve it by least squares --- the standard projection
machinery from linear algebra (Chapter~42). Reading the center back out is completing
the square: $\vec{b} = -\tilde{\mat{Q}}^{-1}\vec{g}$, and the shape is $\tilde{\mat{Q}}$
up to the overall scale, which we fix by demanding corrected magnitude $1$.

\begin{whybox}
Why does fitting an ellipsoid give us a \emph{unique} $\mat{A}$? Strictly, it does
not --- and that subtlety matters. From $\mat{Q}$ we can recover $\mat{A}^{-1}$
only up to an arbitrary rotation, because rotating the whole sensor frame leaves
all magnitudes unchanged and so leaves the ellipsoid identical. In symbols, if
$\mat{A}^{-1}$ works then so does $\mat{R}\mat{A}^{-1}$ for any rotation $\mat{R}$.
The standard fix is to demand the corrected matrix be \emph{symmetric}
positive-definite (the unique symmetric square root of $g^2\mat{Q}$). This pins
down a definite answer that makes the magnitude correct in every direction; the
leftover rotation ambiguity is harmless because it just means ``which way is the
sensor's true $x$ axis'' --- and \emph{that} we fix separately by physically
defining the body frame, or by cross-referencing the accelerometer (see the next
chapter on alignment).
\end{whybox}

\begin{happeningbox}
What is really happening in an ellipsoid fit is a sphere being un-deformed. The
matrix square root of $\mat{Q}$ is doing the work of an inverse SVD: it finds the
three principal axes of your data cloud (the eigenvectors), notes how much each
was stretched (the eigenvalues), un-stretches them back to equal length, and
un-rotates them. Subtracting $\vec{b}$ recenters the cloud. After both steps your
once-lopsided cloud of points lies on a sphere of radius $g$. You can see this
directly by plotting the points before and after the correction.
\end{happeningbox}

\subsection{The simplest version: the bounding box}
\label{ch35:bounding-box}

You do not always need the full ten-parameter fit. The simplest useful method,
common in hobbyist firmware, ignores misalignment and assumes $\mat{A}$ is
\emph{diagonal} (per-axis scale only). Then bias and scale come straight from the
min and max reading on each axis as you tumble the sensor:
\begin{equation}
  b_i = \frac{\max_i + \min_i}{2}, \qquad
  s_i = \frac{\max_i - \min_i}{2}.
  \label{ch35:minmax}
\end{equation}
The bias is the center of the swing; the scale is its half-range. This is the
``hard-iron / soft-iron'' min-max calibration you will see in five-line Arduino
sketches. It captures bias and per-axis scale (the dominant errors) but
\emph{not} misalignment or non-orthogonality. For a hobby wrist controller it is
often good enough; when it is not, you graduate to the full ellipsoid fit.

\begin{example}[Six-position accelerometer calibration by hand]
\label{ch35:ex-sixpos}
A classic deterministic alternative to tumbling is the six-position method. Place the
accelerometer flat in six precise orientations: each of $\pm x$, $\pm y$, $\pm z$
pointing straight down. In each pose, exactly one axis should read $\pm 1\,g$ and
the others $0$. Suppose for the $x$-axis you measure $+1.04\,g$ with $+x$ down and
$-0.98\,g$ with $-x$ down. Solve the two-point line $y = s_x a + b_x$ with
$a=+1g\Rightarrow y=1.04$ and $a=-1g\Rightarrow y=-0.98$:
\[
  s_x = \frac{1.04-(-0.98)}{2} = 1.01, \qquad
  b_x = \frac{1.04+(-0.98)}{2} = 0.03\,g.
\]
So this axis has a $1\%$ scale error and a $30\,\mathrm{mg}$ bias. Corrected
reading: $\hat a_x = (y_x - 0.03)/1.01$. Repeat for $y$ and $z$. Six poses give
you the diagonal-$\mat{A}$ calibration without any curve fitting --- the
geometry of $\pm 1g$ does the work. Cross-axis terms need the ellipsoid fit or
extra poses.
\end{example}

\subsection{Magnetometer: hard iron and soft iron}
\label{ch35:mag-cal}

\begin{historybox}
\textbf{Demoz Gebre-Egziabher and Attitude-Independent Ellipsoid Fitting (2001)}
For nearly a century, calibrating an aircraft or marine magnetic compass required a costly physical procedure. In naval aviation, an entire aircraft had to be towed to a surveyed, non-magnetic compass rose on the tarmac or placed inside an enormous non-magnetic Helmholtz cage, rotating the airframe through precisely known heading angles to record the magnetic deviation curves. As strapdown solid-state magnetometers arrived in the 1990s, the same bottleneck persisted: calibrating the $3\times 3$ soft-iron matrix and $3\times 1$ hard-iron vector required multi-axis rate tables in specialized non-magnetic facilities costing tens of thousands of dollars.

In 2001 at Stanford University, Demoz Gebre-Egziabher, Gabriel Elkaim, J.~David Powell, and Bradford Parkinson solved this with a geometric argument. They used the fact that the Earth's local geomagnetic field vector has a constant magnitude $B_0 = \|\vec{B}_{\text{world}}\|$. If an uncalibrated magnetometer is tumbled through arbitrary 3D orientations, its true measurements must lie on a sphere of radius $B_0$. Hard-iron biases $\vec{b}$ translate this sphere away from the origin; soft-iron materials (nickel battery casings, copper traces, ferrite cores) distort and rotate the sphere into an arbitrary tilted 3D ellipsoid.

Gebre-Egziabher formulated the problem as an attitude-independent \textbf{two-step least-squares ellipsoid fit}. Because the formulation requires no external attitude reference, no turntables, and no optical tracking, a user can calibrate a magnetometer anywhere on Earth simply by tumbling the device. This 2001 method became the standard basis for magnetometer calibration in aerospace, consumer smartphones (the familiar ``wave your phone in a figure-8'' motion), and wearable motion controllers. An open problem is real-time dynamic recalibration as nearby steel objects and fluctuating battery currents warp the field during operation.
\end{historybox}

\begin{trapbox}
\textbf{Why Min--Max Calibration Fails for Soft Iron}
A common beginner shortcut in hobbyist firmware is to apply the per-axis min--max
recipe of \eqref{ch35:minmax} to a magnetometer and stop there. That formula assumes that the distortion matrix $\mat{A}$ is purely \emph{diagonal} and aligned with the chip's coordinate axes. But real soft iron does not merely stretch the axes independently --- it introduces \emph{off-diagonal shear} ($A_{xy}, A_{xz}, A_{yz}$) that tilts the principal axes of the ellipsoid relative to the sensor body.

If soft iron rotates the ellipse by, say, $20^\circ$ relative to the sensor frame, the min--max bounding box is an axis-aligned bounding rectangle that misidentifies the true ellipsoid center and semi-axis lengths. Applying min--max calibration leaves the ellipsoid squashed and rotated, producing heading errors of $15^\circ$ to $30^\circ$ that vary sinusoidally with heading. 
In addition, naive unconstrained algebraic quadric fitting ($\vec{y}^\top \mat{Q}\vec{y} + 2\vec{g}^\top\vec{y} + d = 0$) on incomplete or noisy tumbling data can collapse into an open hyperboloid or cylinder instead of an ellipsoid. Enforcing positive definiteness ($\det \mat{Q} > 0$) via constrained least squares is needed to guarantee a closed, stable ellipsoidal shell.
\end{trapbox}

For magnetometers the ellipsoid distortions have traditional names. \emph{Hard-iron}
distortion is a constant added field from permanent magnets and magnetized metal
fixed to your board (a speaker magnet, a steel screw). It \emph{shifts the sphere
center} --- it is exactly the bias $\vec{b}$. \emph{Soft-iron} distortion comes
from nearby ferromagnetic material that bends and concentrates the field
direction-dependently; it \emph{stretches and tilts the sphere} into an ellipsoid
--- it is exactly the matrix $\mat{A}$. So the ellipsoid model and the traditional
magnetometer terms describe the same thing: hard iron $=\vec{b}$, soft iron
$=\mat{A}$. The cure for both is one ellipsoid fit.

\begin{trapbox}
Soft-iron sources that move \emph{with} the sensor (screws on your own board) can
be calibrated out, because they distort the field consistently in sensor
coordinates. Iron in the \emph{environment} (a steel desk, a passing car, the
motor in a nearby drone) cannot --- it changes as you move through the room, so it
is not a fixed function of your orientation. Magnetometer calibration assumes all
the iron rides along with the sensor. Out in a steel-framed building, even a
perfectly calibrated magnetometer will give wrong headings. This is a main reason a
wrist tracker should not rely heavily on the magnetometer for heading.
\end{trapbox}

\begin{lstlisting}[language=Python, caption={Hard/soft-iron-aware ellipsoid fit (least squares), then correction.\label{ch35:lst-ellipsoid}}]
import numpy as np

def fit_ellipsoid(Y):
    # Y: (N,3) magnetometer/accel samples taken while tumbling the sensor.
    # Fit  x^T Q x + 2 g^T x = 1  (general quadric), then extract center b
    # and shape matrix; return corrector W and offset b so that
    #   a_hat = W @ (y - b)  has constant unit magnitude.
    x, y, z = Y[:, 0], Y[:, 1], Y[:, 2]
    # Design matrix for the 9 quadric coefficients (RHS folded to 1).
    D = np.column_stack([x*x, y*y, z*z, 2*x*y, 2*x*z, 2*y*z, 2*x, 2*y, 2*z])
    rhs = np.ones_like(x)
    coef, *_ = np.linalg.lstsq(D, rhs, rcond=None)    # least-squares solve
    A11, A22, A33, A12, A13, A23, gx, gy, gz = coef
    Q = np.array([[A11, A12, A13],
                  [A12, A22, A23],
                  [A13, A23, A33]])
    # If the origin lies outside the ellipsoid (offset larger than radius),
    # normalizing the RHS to +1 flips the sign of the whole fit. Undo that.
    if np.trace(Q) < 0:
        Q, gx, gy, gz = -Q, -gx, -gy, -gz
    # Completing the square (the 2s in D cancel): center b = -Q^{-1} g.
    b = -np.linalg.solve(Q, np.array([gx, gy, gz]))   # ellipsoid center
    # Shape: symmetric square root of Q, scaled so corrected radius = 1.
    evals, evecs = np.linalg.eigh(Q)
    evals = np.clip(evals, 1e-12, None)               # guard tiny negatives
    Qsqrt = evecs @ np.diag(np.sqrt(evals)) @ evecs.T
    scale = np.mean([np.linalg.norm(Qsqrt @ (yi - b)) for yi in Y])
    W = Qsqrt / scale                                 # so |W (y-b)| ~ 1
    return W, b

# usage:  a_hat = W @ (y - b)   -> lies on the unit sphere
\end{lstlisting}

\begin{labbox}
Log raw magnetometer readings for thirty seconds while you slowly rotate the
sensor in a figure-eight, covering as many orientations as you can, then make a
3D scatter (or three 2D scatters: $xy$, $xz$, $yz$). \emph{Before} calibration
the points form a tilted, off-center ellipsoidal shell. Run
Listing~\ref{ch35:lst-ellipsoid}, apply $\hat{\vec{a}} = \mat{W}(\vec{y}-\vec{b})$,
and scatter again: the shell becomes a unit sphere centered on the origin. Now
plot $\norm{\hat{\vec{a}}}$ for every sample --- it should hug $1$ flat across all
orientations. That flat line is the signature of a good calibration.
\end{labbox}

\section{Gyroscope bias estimation}
\label{ch35:gyro-bias}

The sphere method does not work for a gyroscope. A gyro measures \emph{rate of rotation}, not a constant-magnitude
field. When the device is still, the true rate is zero --- there is no sphere to
trace, just a point at the origin. So gyros need a different, and in some ways
simpler, calibration. The dominant systematic error is bias, and the dominant
problem bias causes is \emph{integration drift}.

\subsection{Why gyro bias dominates long-term error}
\label{ch35:gyro-drift}

To get orientation from a gyro you integrate rate over time: angle is the running
sum of rate $\times$ $\Delta t$. Now suppose the gyro has a tiny constant bias
$b$. Integrating a constant gives a ramp:
\begin{equation}
  \theta_{\text{error}}(t) = \int_0^t b \, \dd\tau = b\,t.
  \label{ch35:bias-integrates}
\end{equation}
A bias of just $0.5^\circ/\mathrm{s}$ --- too small to notice in the raw rate ---
becomes a $30^\circ$ heading error after one minute and a full $180^\circ$ after
six. This is why, of all the error terms, gyro bias is the one you must hunt down
most aggressively. It is also why a gyro alone can never hold orientation
indefinitely: even after you subtract the best constant bias you can measure, the
\emph{leftover} drifting bias (the rate random walk of
Section~\ref{ch35:random-walk}) keeps integrating into an ever-growing angle
error. The accelerometer and magnetometer are in the fusion filter to
correct this drift --- which they can only do if \emph{they} have been
calibrated first.

\begin{figure}[H]
\centering
\begin{tikzpicture}
\begin{axis}[
title style={font=\small},
    width=11.5cm, height=4.8cm,
    xlabel={time (s)}, ylabel={integrated angle error ($^\circ$)},
    xmin=0, xmax=60, ymin=-2, ymax=35,
    axis lines=left, samples=80, domain=0:60,
    legend pos=north west, legend cell align=left,
    title={A constant gyro bias integrates into linear drift},
    every axis plot/.append style={line width=1.0pt},
]
\addplot[specred] {0.5*x}; \addlegendentry{bias 0.5 deg/s ramp}
\addplot[sigblue] {0.1*x}; \addlegendentry{bias 0.1 deg/s, after coarse cal}
\addplot[targetgreen] {0.03*x + 0.9*sin(deg(x/9))};
  \addlegendentry{residual random-walk}
\end{axis}
\end{tikzpicture}
\caption{Integrating rate to get angle turns a constant bias into a linear ramp
(red). Subtracting a measured bias flattens the ramp (blue), but the leftover
slowly drifting bias still wanders (green). No static calibration removes the
green term; that is the job of sensor fusion.}
\label{ch35:fig-gyro-drift}
\end{figure}

\subsection{Estimating the bias: average while still}
\label{ch35:gyro-bias-estimate}

The good news: because the true rate is zero when the device is motionless, gyro
bias is trivial to measure --- \emph{hold the device still and average}. With $N$
samples $\omega_k$ taken at rest,
\begin{equation}
  \hat{b} = \frac{1}{N}\sum_{k=1}^{N}\omega_k.
  \label{ch35:gyro-bias-avg}
\end{equation}
This is Chapter~20's averaging law: the true
signal is a known constant (zero), so the sample mean estimates the bias, and the
white-noise part shrinks like $1/\sqrt{N}$. Average a few seconds at, say,
$200\,\mathrm{Hz}$ and you have a clean bias estimate. The catch, as noted
above, is that $\hat b$ is only good \emph{now}: it drifts with temperature
and time. So good firmware re-runs \eqref{ch35:gyro-bias-avg} automatically every
time it detects stillness.

\begin{whybox}
How does the device know it is still, so it can re-estimate bias? You watch two
cheap signals. First, the accelerometer magnitude: at rest it reads almost
exactly $1\,g$ with very little variance, so $\bigl|\,\norm{\vec{a}} - g\,\bigr|$
small over a short window is strong evidence of stillness. Second, the gyro's own
short-term variance: if the rate is barely wiggling, you are not turning. When
both say ``still'' for, say, half a second, you trust that the true rate is zero
and fold the current average into your running bias estimate. This ``zero-rate
update'' (ZRU, or ZUPT for the velocity version) is one of the most effective
ways to limit gyro drift in a wearable, and it is cheap to compute.
\end{whybox}

\section{Allan variance: noise across time scales}
\label{ch35:allan}

We have claimed twice now that white noise shrinks with longer averaging while
bias drift grows. The Allan variance plot shows both effects at
once and lets you read off the numbers. It was invented by David Allan in 1966 to
characterize atomic-clock stability, and it was later applied to inertial sensors because
the problem is the same: how stable is a thing's output over different time
scales?

\begin{whybox}
\textbf{Why ordinary variance fails for drifting sensors:}
David W. Allan, working at the U.S. National Bureau of Standards on the stability
of atomic frequency standards, faced a problem ordinary variance could not
answer. The standard variance of a drifting signal does not converge --- it just
grows as you record longer, because slow drift keeps adding spread. Allan's
fix (1966) was to measure variance \emph{between adjacent averaging windows} of
length $\tau$, and to plot it as a function of $\tau$. The resulting curve cleanly
separates noise processes that live at different time scales. Decades later, IEEE
standard 952 adopted the Allan variance as the canonical way to specify
fiber-optic and MEMS gyroscopes, which is why every serious IMU datasheet quotes
an ``angle random walk'' and a ``bias instability'' number --- both read straight
off an Allan deviation plot.
\end{whybox}

\subsection{What it measures}
\label{ch35:allan-defn}

Take a long recording of a stationary sensor. Pick an averaging time $\tau$.
Chop the record into consecutive bins of length $\tau$, compute the average value
$\bar{y}_i(\tau)$ in each bin, and then measure how much \emph{adjacent} bin
averages differ from one another:
\begin{equation}
  \sigma_A^2(\tau) = \frac{1}{2}\,\E\!\left[\bigl(\bar{y}_{i+1}(\tau) - \bar{y}_i(\tau)\bigr)^2\right].
  \label{ch35:allan-defn-eq}
\end{equation}
The factor of $\tfrac12$ makes the answer match ordinary variance for white
noise. We sweep $\tau$ from tiny to large and plot the \emph{Allan deviation}
$\sigma_A(\tau)$ (the square root) on log--log axes. Different physical noise
processes show up as straight-line segments with characteristic slopes:

\begin{itemize}
\item \textbf{Angle/velocity random walk (white noise)} appears as a slope of
$-\tfrac12$: doubling the averaging time cuts the deviation by $\sqrt 2$. This is
the averaging-helps regime.
\item \textbf{Bias instability} appears as a \emph{flat} region (slope $0$): the
floor of the curve. Averaging longer than this stops helping --- the bias is
wandering as fast as you can average it out. The minimum of the curve is the
sensor's best achievable stability.
\item \textbf{Rate random walk} appears as a slope of $+\tfrac12$: at long
averaging times the deviation \emph{grows}, because slow drift dominates.
\end{itemize}

\begin{figure}[H]
\centering
\begin{tikzpicture}
\begin{axis}[
title style={font=\small},
    width=12cm, height=6.2cm,
    xlabel={averaging time $\tau$ (s, log scale)},
    ylabel={Allan deviation $\sigma_A(\tau)$ (log)},
    xmode=log, ymode=log,
    xmin=0.01, xmax=1000, ymin=0.002, ymax=1,
    axis lines=left, grid=major, grid style={gray!18},
    title={The classic Allan deviation ``bathtub''},
    every axis plot/.append style={line width=1.1pt},
    samples=60,
]
% white-noise branch: slope -1/2
\addplot[sigblue, domain=0.01:7.44] {0.06*x^(-0.5)};
% rate-random-walk branch: slope +1/2
\addplot[specred, domain=23.9:1000] {0.0045*x^(0.5)};
% flat bias-instability floor
\addplot[targetgreen, domain=7.44:23.9] {0.022};
\node[sigblue, anchor=south west] at (axis cs:0.012,0.45) {slope $-1/2$: white noise (ARW)};
\node[targetgreen, anchor=north, fill=white, inner sep=2pt] at (axis cs:13,0.016) {flat: bias instability};
\node[specred, anchor=south east] at (axis cs:900,0.10) {slope $+1/2$: rate random walk};
\end{axis}
\end{tikzpicture}
\caption{The Allan deviation of a real MEMS sensor traces a ``bathtub.'' At short
averaging times white noise dominates and falls with slope $-1/2$ (averaging
helps). It bottoms out at the bias-instability floor (flat). At long times slow
drift takes over and the curve climbs with slope $+1/2$. The bottom of the
bathtub tells you the best averaging time and the best stability you can hope
for.}
\label{ch35:fig-allan}
\end{figure}

\begin{intuitionbox}
Read an Allan deviation curve from left to right across three time-scale regimes.
At short averaging times ($\tau \ll 1\,\mathrm{s}$), broadband white noise dominates and
averaging drives the uncertainty down as $\tau^{-1/2}$. Around a few seconds, the curve
bottoms out at its bias-instability floor --- the lowest uncertainty the sensor can achieve.
Average any longer and slow random-walk bias drift takes over, pushing the error back up as
$\tau^{+1/2}$. The $\tau$ at the minimum of that bathtub curve is the optimal window length for
stationary zero-rate updates: shorter windows leave white noise unaveraged, while longer windows
let bias wander contaminate the estimate.
\end{intuitionbox}

\begin{connectionbox}
The Allan variance is a cousin of the power spectral density from Chapter~20.
Both describe how a random signal's energy is distributed across time scales /
frequencies; in fact each Allan-slope corresponds to a power-law region of the
PSD ($1/f^0$ white, $1/f$ flicker, $1/f^2$ random walk). The Allan variance is
often preferred for sensors because the integrating, drifting low-frequency
behavior that ruins a gyro is exactly where the ordinary PSD estimate is hardest
to pin down --- and where the Allan plot is clearest and most stable.
\end{connectionbox}

\begin{lstlisting}[language=Python, caption={Computing the Allan deviation from a long stationary recording.\label{ch35:lst-allan}}]
import numpy as np

def allan_deviation(y, fs, taus):
    # y: 1D array of a stationary sensor axis; fs: sample rate (Hz)
    # taus: list of averaging times (s) to evaluate
    out = []
    for tau in taus:
        m = int(tau * fs)                 # samples per averaging bin
        if m < 1 or 2*m > len(y):
            continue
        n = len(y) // m
        bins = y[:n*m].reshape(n, m).mean(axis=1)   # bin averages
        d = np.diff(bins)                            # adjacent differences
        sigma2 = 0.5 * np.mean(d**2)                 # Allan variance
        out.append((tau, np.sqrt(sigma2)))
    return np.array(out)
\end{lstlisting}

\section{Why calibration must come before fusion}
\label{ch35:before-fusion}

This section states the chapter's main claim. The
next chapter builds a \emph{sensor fusion} filter that blends gyro, accelerometer,
and magnetometer into one orientation estimate. The blend works by \emph{trusting
each sensor where it is strong and correcting it where it is weak}: the gyro for
fast, short-term rotation; the accelerometer and magnetometer to pin down the
slow drift. That correction is only as good as the sensors' agreement about
``which way is down'' and ``which way is north.'' If the accelerometer's down is
tilted $3^\circ$ by uncorrected misalignment, or the magnetometer's north is bent
$20^\circ$ by uncorrected hard iron, the fusion filter will pull your
orientation \emph{toward the wrong answer}, because it has no way to know the
reference is wrong.

\begin{aestheticbox}
Calibration is upstream of everything. A fusion filter, a gesture classifier, a
machine-learning model --- each assumes its inputs mean what they claim. Bad inputs
here produce \emph{plausible-looking} wrong outputs, which is worse than obviously
wrong output because it passes a casual check. The half-hour you spend calibrating
improves every stage downstream. It is a cheap step with a large payoff, and
beginners often skip it. Do it, and do it at operating temperature.
\end{aestheticbox}

\begin{goalbox}
For the wrist controller, here is the calibration procedure for each
sensor. (1) \textbf{Accelerometer:} full
ellipsoid fit once per build, so ``down'' is trustworthy to a fraction of a
degree --- this anchors pitch and roll. (2) \textbf{Magnetometer:} ellipsoid fit
\emph{with the final enclosure, battery, and any magnets in place}, because the
hard/soft iron you must cancel is the iron that ships on the wrist; and accept
that environmental iron will still corrupt heading, so the design must
not depend on absolute magnetic north for playing notes. (3) \textbf{Gyro:}
continuous zero-rate bias updates whenever the wrist is detected still between
gestures --- this is what keeps fast wrist flicks from accumulating drift over a
song. The same approach carries over to the band's tendon and contact sensors,
the channels that actually carry finger movement: the force-myography (FMG) array
and the high-rate accelerometer used as a contact/bio-acoustic sensor each have
their own baseline and gain, and each gets calibrated during the guided
enrollment of the capstone --- the user performs every finger movement a few
times, the phone records the per-finger signal signature, and it sets the matching
thresholds. With that per-finger calibration in hand, finger movements \emph{are}
detectable from a wrist or forearm band: tendon displacement (FMG) and the surface
EMG pulse both carry which finger moved
(Figure~\ref{ch35:fig-fmg-trace}).

One limit remains: all of the IMU work above calibrates the
\emph{wrist's} bulk motion, and continuous per-joint finger angles are simply not
present in that bulk motion --- the problem is underdetermined, and no fit recovers
what the signal does not contain. A continuous finger-pose curve needs sensors on
the fingers themselves (instrumented rings). Treat that as an upgrade path, not a
prerequisite: detecting which finger moved, and when, is enough to play notes, and
that is achievable.
\end{goalbox}

\begin{figure}[H]
\centering
\begin{tikzpicture}
\begin{axis}[
title style={font=\small},
    width=12cm, height=5.2cm,
    xlabel={time (s)}, ylabel={FMG channel (calibrated)},
    xmin=0, xmax=6, ymin=-0.1, ymax=1.15,
    axis lines=left, samples=160, domain=0:6,
    legend pos=north east, legend cell align=left,
    title={Calibrated tendon (FMG) signal: a finger press above its threshold},
    every axis plot/.append style={line width=1.0pt},
]
% baseline subtracted, gain normalized; a press is a bump
\addplot[sigblue, domain=0:6]
  {0.03*sin(deg(3*x)) + 0.85*exp(-((x-3)^2)/0.18)};
\addlegendentry{index-finger channel}
\addplot[trapred, dashed, domain=0:6] {0.4};
\addlegendentry{enrollment threshold}
\node[anchor=south, sigblue] at (axis cs:3,0.9) {\small press detected};
\end{axis}
\end{tikzpicture}
\caption{What per-finger calibration produces on a tendon (FMG) channel. The
enrollment step subtracts each channel's resting baseline and normalizes its gain,
so a still hand reads near zero. When the calibrated index finger presses, that
channel rises into a clear bump; crossing the threshold set during enrollment (red
dashed) fires a per-finger event. No machine learning is involved --- the detector
is a calibrated baseline, a gain, and a threshold per finger, matched live against
the enrolled signature.}
\label{ch35:fig-fmg-trace}
\end{figure}

\begin{buildbox}
A practical calibration routine you can ship in firmware and run on the bench.

\textbf{Gyro bias (every boot, plus continuously).}
\begin{enumerate}
\item At boot, ask the user to set the device down. Confirm stillness:
$\bigl|\,\norm{\vec{a}}-g\,\bigr| < 0.05\,g$ and gyro variance below a threshold
over a $0.5\,\mathrm{s}$ window.
\item Average the gyro over $1$--$2\,\mathrm{s}$ (Eq.~\ref{ch35:gyro-bias-avg})
and store $\hat{\vec b}_{\text{gyro}}$.
\item During use, re-run the average whenever stillness is detected; blend the new
estimate into the stored one with a slow low-pass so temperature drift is tracked.
\end{enumerate}

\textbf{Accelerometer (once per build, at operating temperature).}
\begin{enumerate}
\item Either do the six-position method (Example~\ref{ch35:ex-sixpos}) for the
diagonal model, or tumble slowly through all orientations and run the ellipsoid
fit (Listing~\ref{ch35:lst-ellipsoid}).
\item Store $\mat{W}_{\text{acc}}$ and $\vec{b}_{\text{acc}}$; apply
$\hat{\vec a}=\mat{W}_{\text{acc}}(\vec y-\vec b_{\text{acc}})$ on every sample.
\item Sanity check: corrected $\norm{\hat{\vec a}}$ should be $1\,g\pm$ a few
$\mathrm{mg}$ in every pose.
\end{enumerate}

\textbf{Magnetometer (once, with final hardware assembled).}
\begin{enumerate}
\item Power the board exactly as it will ship (battery in, enclosure closed).
\item Rotate in slow figure-eights for $20$--$30\,\mathrm{s}$, logging raw mag.
\item Run the ellipsoid fit; store $\mat{W}_{\text{mag}}$ (soft iron) and
$\vec{b}_{\text{mag}}$ (hard iron).
\item Verify corrected $\norm{\hat{\vec m}}$ is flat across orientations.
\end{enumerate}

\textbf{Characterize (once, informative).} Log a stationary gyro for $20$--$60$
minutes and compute the Allan deviation (Listing~\ref{ch35:lst-allan}). The
bottom-of-bathtub $\tau$ tells you how long to average for zero-rate updates; the
floor tells you the best heading stability to expect from the gyro alone.
\end{buildbox}

\begin{applicationbox}
This exact ellipsoid machinery shows up far beyond wearables. Smartphone compass
apps run a continuous background magnetometer ellipsoid fit --- that ``wave your
phone in a figure eight'' prompt is collecting points to refit the sphere after
hard-iron conditions change. Drone flight controllers calibrate accelerometer
and magnetometer with the same six-position and tumble routines before every
serious flight, because a $2^\circ$ heading error sends an autonomous waypoint
mission off course. Even MRI and geophysical magnetometers use sphere/ellipsoid
calibration to remove fixed instrument bias. The geometry of ``constant magnitude
implies a sphere'' is widely reused in sensing.
\end{applicationbox}

\section{Worked example: separating bias from drift}
\label{ch35:worked}

\begin{example}[Reading an Allan plot for averaging time]
\label{ch35:ex-allan}
Suppose your gyro's Allan deviation reaches its minimum (flat floor) at
$\tau^\star \approx 8\,\mathrm{s}$ with a floor value of $0.02^\circ/\mathrm{s}$,
and at short times the white-noise branch passes through
$\sigma_A(1\,\mathrm{s}) = 0.06^\circ/\mathrm{s}$ with slope $-1/2$. Two practical
readings follow. First, if you average a zero-rate update over $8\,\mathrm{s}$ you
get the best possible bias estimate, good to about $0.02^\circ/\mathrm{s}$;
averaging longer is wasted effort because the floor is flat and then climbs.
Second, the angle random walk (the white-noise branch's value at
$\tau=1\,\mathrm{s}$, by the standard convention) is
$N \approx 0.06^\circ/\sqrt{\mathrm{s}}$, so after $t$ seconds of free
integration the random-walk angle uncertainty is roughly
$0.06\sqrt{t}\ ^\circ$ --- about $0.4^\circ$ after $40\,\mathrm{s}$. That tells
you how often the accelerometer must step in to reset drift, which is exactly the
tuning knob of the next chapter's fusion filter.
\end{example}

\begin{example}[Bias still integrates faster than random walk]
\label{ch35:ex-bias-vs-walk}
Contrast Example~\ref{ch35:ex-allan} with an \emph{uncalibrated} bias of
$0.5^\circ/\mathrm{s}$. After $40\,\mathrm{s}$ that bias alone contributes
$0.5\times40 = 20^\circ$ of error --- fifty times the $\sim 0.4^\circ$ from random
walk. A constant bias you failed to remove is far larger than the noise floor.
Removing the deterministic $20^\circ$ is what makes the device usable; noise
performance matters only after the systematic part is gone.
\end{example}

\section{Sensor Models in Practice: Calibration Bounds, Real-World Trade-Offs, and the Path to Fusion}
\label{ch35:capabilities-limits}

\subsection{What Calibration Achieves}
Calibration turns distorted raw readings into accurate physical vectors:
\begin{itemize}
  \item \emph{In-situ user calibration:} By using least-squares ellipsoid fitting, a user waving a smartwatch or smartphone in a figure-eight can calibrate out onboard battery and speaker distortions in seconds without specialized test fixtures.
  \item \emph{Sub-degree attitude accuracy:} Extracting deterministic bias $\vec{b}$ and whitening matrix $\mat{W} = \mat{A}^{-1}$ reduces accelerometer tilt errors from several degrees to under $0.1^\circ$, and cuts magnetometer heading errors from over $30^\circ$ down to $<1^\circ$ in magnetically clean environments.
  \item \emph{Noise floor awareness:} The Allan variance plot provides a quantitative boundary separating high-frequency white noise (which can be averaged down) from low-frequency bias instability (where averaging actively degrades the estimate), establishing the optimal fusion update rate.
\end{itemize}

\subsection{Limitations}
Calibration solves only \emph{static, repeatable} errors. It cannot fix dynamic or environmental shifts:
\begin{itemize}
  \item \emph{Environmental magnetic distortion:} Calibrating out the soft iron of your wristband's battery case does nothing when you walk near a reinforced concrete floor, pass through a steel subway turnstile, or place your hand near a laptop keyboard. The external magnetic field vector itself changes, so even a calibrated magnetometer reports false yaw.
  \item \emph{Thermal drift:} Semiconductor biases drift with temperature ($0.1\,\text{mg}/^\circ\text{C}$ for accelerometers; $0.05^\circ/\text{s}/^\circ\text{C}$ for gyros). A calibration performed on a cool morning will drift when your wrist warms the PCB during a performance.
  \item \emph{Inertial motion contamination:} An accelerometer calibrated to $1\,g$ still cannot distinguish the static Earth gravity vector from dynamic arm swings during rapid gestures.
\end{itemize}

\subsection{Next: From Calibration to Multi-Sensor Fusion}
Calibration gives us clean, unbiased sensor measurements, but each individual sensor still has errors over time: the gyroscope rotates cleanly but drifts indefinitely, while the accelerometer and magnetometer provide absolute references that are corrupted by dynamic motion and magnetic anomalies. 

In \textbf{Chapter 36}, we bring all three sensors together in \textbf{Sensor Fusion}. We will use quaternion mathematics to avoid gimbal lock, and develop complementary, Kalman, and Madgwick filters to fuse high-frequency gyro responsiveness with low-frequency gravitational and magnetic stability, to obtain drift-free, real-time 3D orientation tracking for our gesture instrument.

\section{Exercises}
\label{ch35:exercises}

\begin{exercise}
A flat accelerometer's $z$-axis reads $+1.06\,g$ pointing up and $-0.96\,g$
pointing down. Find the scale factor $s_z$ and bias $b_z$ for the affine model
$y=s_z a + b_z$, and write the correction formula for $\hat a_z$. Is the bias
positive or negative?
\end{exercise}

\begin{exercise}
Explain, using the input--output line of Figure~\ref{ch35:fig-bias-scale}, why a
single measurement at $+1\,g$ cannot distinguish bias from scale error. What is
the minimum number of distinct true input levels needed to separate them on one
axis, and why?
\end{exercise}

\begin{exercise}
A magnetometer's raw data, plotted in 3D while tumbling, forms a sphere whose
center is at $(12, -8, 30)\ \mu\mathrm{T}$ rather than the origin. Which error
term is this --- hard iron or soft iron? What is the corresponding vector in the
model of \eqref{ch35:lumped-model}, and how do you correct it?
\end{exercise}

\begin{exercise}
Using the min--max method of \eqref{ch35:minmax}, suppose the $x$-axis of a
magnetometer swings between $\min_x = -55$ and $\max_x = +5$. Compute the
hard-iron offset and per-axis scale for $x$. If the $y$-axis swings between
$-30$ and $+30$, what does the difference in scale tell you about soft iron?
\end{exercise}

\begin{exercise}
A gyro has a constant bias of $0.2^\circ/\mathrm{s}$ that you did not remove. You
integrate to track heading. How long until the heading error reaches $90^\circ$?
Now suppose you remove that bias but a residual \emph{angle random walk} makes the angle
uncertainty grow as $0.05\sqrt{t}\ ^\circ$. How long until \emph{that} reaches
$90^\circ$? Comment on the enormous ratio.
\end{exercise}

\begin{exercise}
Sketch (by hand or in code) the Allan deviation bathtub. Label the slope of each
of the three regions and name the noise process responsible. Mark the point that
gives the optimal averaging time, and explain in one sentence why averaging
longer than that point makes the estimate \emph{worse}.
\end{exercise}

\begin{exercise}
Why does the sphere-fit calibration work for accelerometers and magnetometers but
\emph{not} for gyroscopes? Identify the geometric property the first two share
that the gyro lacks, and describe the alternative calibration the gyro requires
instead.
\end{exercise}

\begin{exercise}
You calibrate a magnetometer perfectly on your wooden workbench, then carry the
device to a steel-framed doorway. The calibration that gave a flat
$\norm{\hat{\vec m}}=1$ on the bench now gives readings off by $15^\circ$ in
heading. Explain why, distinguishing iron that moves with the sensor from iron
fixed in the environment. What does this imply for a wrist controller's use of
magnetic heading?
\end{exercise}

\begin{exercise}
Write pseudocode for a stillness detector suitable for triggering a zero-rate
gyro bias update. Use both the accelerometer magnitude and the short-window gyro
variance. State the two thresholds you would tune and what failure mode each one
guards against.
\end{exercise}

\begin{exercise}[Observability and model limits]
A common calibration pitfall is confusing \emph{parameter identification} with
\emph{state observability}. Explain why calibrating the $3\times 3$ scale/misalignment
matrix $\mat{A}$ and $3\times 1$ bias $\vec{b}$ of a wrist-worn IMU eliminates sensor-frame
coordinate distortion, but cannot expand the dimensionality of what the sensor measures.
Specifically, contrast the rigid-body degree-of-freedom tracking of the forearm with
the high-dimensional internal state space of finger tendon articulation, stating what
additional physical sensing channel is physically required to observe tendon excursion.
\end{exercise}

